\documentclass{article}
\usepackage[margin=2cm]{geometry}
\usepackage{amsmath}
\usepackage{graphicx}
\usepackage{booktabs}
\usepackage[utf8]{inputenc}
\usepackage{ragged2e}
\setlength{\emergencystretch}{3em}
\begin{document}
\sloppy

\subsection*{1 Gamma matrix and spinor conventions}


ncreteness, we take the following basis of gamma matrices


\begin{equation}
\begin{array} { l } { \gamma _ { 1 } = \sigma _ { 2 } \otimes \mathbf { 1 } _ { 2 } \otimes \sigma _ { 3 } } \\ { \gamma _ { 2 } = \sigma _ { 2 } \otimes \mathbf { 1 } _ { 2 } \otimes \sigma _ { 1 } } \\ { \gamma _ { 3 } = 1 _ { 2 } \otimes \sigma _ { 1 } \otimes \sigma _ { 2 } } \\ { \gamma _ { 4 } = 1 _ { 2 } \otimes \sigma _ { 3 } \otimes \sigma _ { 2 } } \\ { \gamma _ { 5 } = 1 _ { 1 } \otimes \sigma _ { 3 } \otimes \sigma _ { 2 } } \\ { \gamma _ { 6 } = \sigma _}\end{array}
\end{equation}


These gamma matrices satisfy the Clifford algebra


\begin{equation}
\begin{array} { r } { \left \{ \gamma _ { \mu }, \gamma _ { \nu } \right \} = 2 \delta _ { \mu \nu } } \\ { \mathrm { ~ a ~ l u t h ~ o u s. } [ \mathrm { ~ b e l u c h ~ i n t h ~ o u s h ~ d i n t h ~ o u s h ~ o u s h ~ o u s h ~ o u s h ~ } _ { \mathrm { ~ m o t h t h ~ o u s h ~ o u s h ~ o u s h ~ o u s h ~ o u s h ~ o u s h ~ o u s h ~ o u s h ~ o u s}}}\end{array}
\end{equation}


as appropriate for a positive definite Euclidean spacetime. All matrices are purely imaginary
and satisfy


\begin{equation}
( \gamma _ { \mu } ) ^ { \dag } = \gamma _ { \mu } ~
\end{equation}


We will now be interested in a seven-dimensional Clifford algebra, which will require the
introduction of a new matrix /7. The reason we are interested in this is that we would like
This allows us to determine properties of the Dirac spinors in the Euclidean-continued F(4)
gauged supergravity theory with H6 background by first considering Dirac spinors in seven
dimensions and then performing a timelike reduction. In particular, we will choose a 7D
metric of signature (+, +, +, +, +, +, -) for the ambient space. Then hyperbolic space H6
is given by the following quadratic form


\begin{equation}
\begin{array} { r l } & { \quad \quad ( 4 ) \qquad x _ { 1 } ^ { 2 } + \dots + x _ { 6 } ^ { 2 } - x _ { 7 } ^ { 2 } = - L ^ { 2 } } \\ & { \quad \vdots, \dots, \dots, \dots, \dots, \dots, \dots, \dots, \dots, \dots, \dots, \dots, \dots, \dots, \dots, \dots, \dots, \dots } \end{array}
\end{equation}


The seven-dimensional Clifford algebra is made up of the set of matrices \{71,..., 76, \textasciitilde{}7\},
with 7 satisfying


\begin{equation}
\begin{array} { r l } & { \left \{ \gamma _ { \mu }, \gamma _ { \Gamma } \right \} = 0 ~ ~ \forall \mu \neq 7 } \\ & { \quad \left \{ \gamma _ { \mu }, \gamma _ { \Gamma } \right \} = 0 ~ ~ \forall \mu \neq 7 } \\ & { \qquad }\end{array}
\end{equation}


As usual, we use the notation  = (7)-', so that by the above we have  = /7. We now
discuss Dirac spinors in d = 7. We define the Dirac conjugate of A to be


\begin{equation}
\begin{array} { c c c c c c c c } { } & { } & { } & { } & { } & { } & { } & { } & { } & { } & { } & { } & { } \\ { } & { } & { } & { \ddots } & { } & { } & { } & { } & { } & { } & { } \\ { } & { } & { } & { } & { } & { } & { } & { } & { } & { } & { ( 6 ) } \\ { } & { } & { } & { } & { } & { } & { } & { } & { } & { }\end{array}
\end{equation}


for some matrix G. There are two possible choices for G , which in the particular case of the
ambient space above are


\begin{equation}
G _ { 1 } = \gamma ^ { 7 } \qquad \qquad \qquad G _ { 2 } = \gamma ^ { 1 } \dots \gamma ^ { 6 } \qquad ( 7 )
\end{equation}

\end{document}
